\section{Methods detail}
\label{supp:methods}

This section carries the methodological detail compressed out of Sections~3,
4 and~5 of the manuscript: the data-construction arguments behind the
structural-zero and admissibility rules, the full specification of each of the
ten model families and each of the seven uncertainty mechanisms, the crossed
grid's design floors and compute arithmetic, the single fit-and-sample
interface every cell implements, the seeding and provenance scheme, and the
estimator algebra behind the metrics.

Nothing here is a result. Everything here is the answer to ``how exactly was
that computed?'' for a number reported in the manuscript, and it is written so
that each subsection can be read on its own. It supports the manuscript's
central methodological premise---that one evaluation code path scores every
cell of the grid, so that a coverage difference between two cells is a
difference between a family or a mechanism and not between two metric
implementations---and it supplies the specifications the manuscript points to
by name: the hyperparameter grids and architectures of
Section~4, EnbPI's four documented deviations, the sorted-sample CRPS
estimator, and the randomised PIT of \citet{czado2009predictive}. The
registered design itself, its addenda and the validation gates are
Section~\ref{supp:prereg}; the full per-cell evidence is Supplementary
Section~S3.

\subsection{Regimes, origins and the splitting discipline}
\label{supp:methods-regimes}

Three regimes are scored, all expanding-origin. The \emph{stable} control
steps its origin two years at a time, $T=1990,1992,\dots,2014$, thirteen
origins in all, each testing $T+1,\dots,T+5$ with every test year at or before
2019, so that no break falls anywhere in any of its test windows: it measures
calibration where nothing happens. The \emph{shift} regime trains to 2019 and
tests 2020--2024 at $h=1,\dots,5$. The \emph{placebo} regime trains to 1913
and tests 1914--1922 at $h=1,\dots,9$, a second event rather than more draws
from the first.

A random split is never used anywhere. Mortality panels leak trivially under
random splitting, because adjacent ages and adjacent years are near-duplicates
of one another: a model that has seen 2019 and 2021 can interpolate 2020
without forecasting anything. Training always uses the maximal contiguous
block of complete years ending at the origin, from the population's own first
year of data, and an (origin, population, sex) triple is admissible only with
at least fifteen such years (Addendum~3~\S2 and~\S4). Hyperparameters are
tuned on an inner \emph{time} split---the final five pre-cutoff years of each
training window---never on a test year, never on a random subset, and never
against coverage, which is the audited quantity.

\subsection{Data construction}
\label{supp:methods-data}

\paragraph{The bulk files.} The bulk deaths and exposures files of the pinned
vintage each carry fifty distinct HMD population codes (the
\texttt{PopName} count recorded in the manifest header): forty-one countries
or areas, plus nine further codes for the overlapping sub-population and
civilian variants HMD publishes for Germany (two), France (one), the United
Kingdom (four) and New Zealand (two). Those nine are the reason the placebo
panel needs an explicit exclusion rule rather than a country list: FRACNP and
GBRCENW are civilian-population variants overlapping FRATNP and GBRTENW, and
keeping both members of a pair would double-count one country.

\paragraph{Structural zeros: why the weighted fit is the registered fit.}
\label{supp:methods-zeros}
A zero-exposure cell inside ages 0--99 is a structural zero, not a missing
value: HMD's 1~January population count is zero at both corners of the Lexis
square, so nobody was alive at that age, and the observed death count is zero
as well. Under the Poisson likelihood with log-exposure offset such a cell has
mean $\lambda = E\exp(\eta) = 0$ and observation $D=0$, so its log-likelihood
contribution, its score and its Fisher information are all identically zero.
The maximum-likelihood estimate computed with the cell included therefore
equals the estimate computed under indicator weights \emph{exactly}, for any
parameter value; the only thing that fails is the numerical evaluation of
$0\log 0$. This is why carrying $W_{x,t}=\ind\{E_{x,t}>0\}$ is a numerically
correct evaluation of the registered objective and not a change of objective,
and why Addendum~3~\S1 records it as a clarification rather than a deviation.

The weighting is implemented per family, and the addendum fixes each path:
weighted Poisson likelihood for Poisson Lee--Carter and Renshaw--Haberman,
with zero-weight cells contributing nothing to the Newton or IWLS steps; a
weighted rank-one fit by alternating least squares or EM-SVD for the SVD
Lee--Carter, with the exact unweighted SVD kept as a fast path when all
weights are one; per-year weighted least squares on the logit of $q$ for CBD,
replacing an unweighted column-mean shortcut for $\kappa^{(1)}$; and, for the
sparse VAR, a common-complete-subsample rule under which each fit uses
regression rows built only from consecutive year triples in which every age of
the banded window is observed. The last is the only clause that loses data,
and the loss was measured before the sweeps (training years to 2019): CHE
female $142\to139$, CHE male $142\to125$, FIN male $140\to101$, ISL female
$180\to164$, ISL male $180\to136$, LUX male $58\to49$, SWE male $267\to242$
regression rows, all far above the banded VAR's minimum, every defective year
before 1970. No positive-semidefiniteness projection is applied to the
residual covariance, because a projection would perturb predictive width---the
audited quantity---by an unquantified amount. The reduction requirement is
tested rather than asserted: every weighted path reproduces the unweighted fit
\emph{bit-identically} on panels with no zero-exposure cells.

\paragraph{Belgium, and why contiguity is a rule.} Belgium's 1914--1918 deaths
and exposures are missing at every age. In the study's modelled unit that is
1{,}000 cells---100 ages $\times$ 2 sexes $\times$ 5 years---and in the raw
$1\times1$ file it is 555 rows: 111 ages including the $110{+}$ group, one row
per age-year, five years, present but carrying HMD's ``\texttt{.}'' missing
marker in the Female, Male and Total columns. Without the contiguity rule the
pivot silently deletes the five missing year columns and the panel closes up,
with two specific consequences: a random walk with drift then treats 1913 to
1919 as a single one-year step, and the Renshaw--Haberman cohort index, which
is positional, mislabels every cohort after 1913 by five years. The rule---the
training window is the maximal contiguous block of complete years ending at
the origin---is what prevents both, and it binds exactly once in the shift
regime, on BEL, which trains on 1919--2019.

\paragraph{The fifteen-year admissibility floor.} Very short training windows
do not fail loudly; they produce plausible-looking junk. A random-walk
innovation variance estimated from two differences is a number, and it will be
written into a coverage column as though it were a forecast: the measured
instance in Addendum~3~\S4 is CHL males at origin 1994, three training years,
which wrote an LC/native 95\% coverage of 0.784 as a valid row. The floor
removes such rows from the panel rather than from the reporting. Its
consequence is that the stable panel is unbalanced across origins, the number
of buildable population--sex series growing as the late-starting Baltic, East
Asian and South American series pass the floor, which is why every
stable-regime summary is read beside its per-origin population count: the wild
cluster bootstrap's behaviour depends on the number of clusters actually
present.

\paragraph{Losses by family.} The likelihood is Poisson with a log-exposure
offset throughout \citep{brouhns2002poisson}, with two deliberate exceptions
retained as historical baselines---the original singular-value Lee--Carter
with its least-squares fit on log rates, and CBD with its per-year
least-squares fit on the logit of $q_{x,t}$. The neural families fit the same
Poisson deviance with the log-exposure offset (neural-LC and the shallow CNN)
or a negative-binomial log-likelihood (the distributional head); the
LSTM-on-$\kappa_t$ inherits the SVD Lee--Carter stage and fits the index by
squared error, as its reference paper does; and the multi-output Gaussian
process places a Gaussian likelihood on the half-count log rate with a
Poisson-level noise floor. Every loss is written in the family specification
\texttt{docs/NEURAL-SPEC.md}, committed before the neural sweep.

\paragraph{Life-table conventions and the 0.15-year tolerance.}
\label{supp:methods-lifetable}
One period-life-table routine serves every predictive sample and the observed
side alike: $q_x = m_x/(1+(1-a_x)m_x)$; $a_0 = 0.07+1.7\,m_0$ capped to
$[0.01,0.35]$, a linear infant-separation rule in the spirit of
\citet{andreev2015average}; $a_x=\tfrac12$ elsewhere; and constant-hazard
closure at the last age, $e_A = 1/m_A$. HMD's own life table instead uses the
piecewise Andreev--Kingkade coefficients for $a_0$ and its own old-age
treatment, so the two tables are not identical by construction. That
difference is what validation gate~3's tolerance measures: our $e_0$ and
$e_{65}$ on HMD's published $m_x$ must reproduce HMD's published values within
0.15~years, which is the size of the documented $a_x$ simplifications rather
than an arbitrary allowance. Because the realised and predictive sides are put
through the same routine, any life-table approximation cancels in every
comparison the paper reports.

\subsection{Model families in full}
\label{supp:methods-families}

Each family is specified by its reference paper and fitted per population and
sex on a matrix of deaths and exposures over ages 0--99 (CBD: 55--99) and
years up to the origin, under the cell weights $W=\ind\{E>0\}$ above.

\begin{enumerate}
\item \textbf{Lee--Carter (SVD).} $\log m_{x,t}=\alpha_x+\beta_x\kappa_t$, a
  rank-one fit to centred log rates under $\sum_x\beta_x=1$ and
  $\sum_t\kappa_t=0$: an exact SVD when no cell is weighted out, an
  expectation--maximisation SVD otherwise, the two agreeing to machine
  precision at unit weights. Forecast by a random walk with drift on
  $\kappa_t$, including drift-estimation uncertainty \citep{lee1992modeling}.
\item \textbf{Poisson Lee--Carter.} The same structure under a Poisson
  likelihood with log-exposure offset, fitted by Newton's method
  \citep{brouhns2002poisson}.
\item \textbf{CBD (M5).} $\operatorname{logit} q_{x,t} = \kappa^{(1)}_t +
  \kappa^{(2)}_t (x-\bar x)$, fitted by per-year weighted least squares on
  ages 55--99, where the logit is near-linear in age, and forecast by a
  bivariate random walk with drift \citep{cairns2006two}. CBD is undefined
  below age~55 and its samples are masked there: it is scored on 45 ages, sits
  in its own row block with that age support stated, and is never averaged
  into a ranking with families scored on 100 ages. Its ``25--64'' column is
  therefore a 55--64 stub.
\item \textbf{APC / Renshaw--Haberman (M2-A).} Lee--Carter with a cohort term
  $\gamma_{t-x}$ under the identifiability constraints of
  \citet{cairns2009quantitative}, with an ARIMA projection of the cohort index
  \citep{renshaw2006cohort}.
\item \textbf{Sparse VAR.} A VAR(1) on age-specific mortality improvements
  $\Delta\log m_{x,t}$ in which age $x$'s equation regresses on the lagged
  improvements of ages within a band of half-width~3, estimated by ordinary
  least squares; the hard band is the sparsity structure, a documented
  simplification of the penalised estimator of \citet{li2017coherent}.
  Innovations are drawn from the cross-age residual covariance shrunk toward
  its diagonal by a fixed weight of $0.2$, and coefficient draws carry the OLS
  coefficient covariance. Regression rows use only year triples in which every
  age of the band is observed (Addendum~3~\S1), and explosive coefficient
  draws are rejected and redrawn rather than clipped (Addendum~3~\S7).
\item \textbf{Multi-output Gaussian process.} The multi-output GP of
  \citet{huynh2021multi} reduced to one population: an exact multitask GP with
  the scaled calendar year as input and ages as tasks, kernel
  $\mathrm{RBF}(\text{year})\otimes\mathrm{ICM}(\text{rank }5)$, a Gaussian
  likelihood on the half-count log rate with a Poisson-level noise floor, and
  training by exact marginal likelihood on the trailing block of complete
  years, the Kronecker structure requiring a complete year-by-age block. Its
  native predictive is the joint posterior over the $h$ future years. The
  window is at least forty years and---by Addendum~4---at most sixty. The
  lower bound is an identification constraint rather than a convention: the
  separable age $\times$ time kernel is fitted by exact marginal likelihood,
  and on a short panel the time-kernel length-scale is not identified against
  the nugget, so the fit can trade smoothness for noise without changing the
  likelihood materially. That is the floor which removes CHL, HKG, HRV and KOR
  from every GP cell.
\item \textbf{Neural-LC.} The embedding network of \citet{richman2021neural}
  in per-panel form: an age embedding of dimension~5 concatenated with the
  scaled year, two hidden layers of 64 tanh units with dropout $0.05$, and a
  log-rate output added to the log-exposure offset under the Poisson deviance.
  The forecast evaluates the network at future years; the extrapolation of the
  year feature is the audited fragility, not a defect to engineer away.
\item \textbf{Shallow CNN-LC.} The convolutional generalisation of
  Lee--Carter \citep{perla2021time}: a trailing window of ten years of
  half-count log rates over all ages, one convolutional layer of eight
  $3\times3$ filters, and a linear map to next year's log rates under the
  Poisson deviance, with recursive multi-step forecasts.
\item \textbf{LSTM-on-$\kappa_t$.} The SVD Lee--Carter age structure with a
  one-layer LSTM of 16 hidden units and input window 10 in place of the random
  walk on $\kappa_t$ \citep{nigri2019deep}. The innovation variance is
  estimated from one-step training residuals, and forecast paths add a fresh
  Gaussian innovation per horizon, exactly as the random walk does.
\item \textbf{Distributional NB head.} The neural-LC features with two outputs
  per cell, $\log\mu$ (plus the offset) and a log-dispersion, trained by the
  negative-binomial (NB2) log-likelihood: the only neural family with a
  model-native predictive distribution. Because the negative binomial is a
  Poisson--gamma mixture, its rate samples are gamma draws, so the Poisson
  composition of Section~\ref{supp:methods-interface} yields the NB count law
  exactly, with no double counting.
\end{enumerate}

\paragraph{Hyperparameter grids.} The neural and GP families tune a learning
rate and an epoch count from grids of two by two values fixed in
\texttt{docs/NEURAL-SPEC.md} before the sweeps, on the inner time split of
Section~\ref{supp:methods-regimes}, by the family's own loss. One grid was
corrected before any real-data run: the CNN's learning rates were lowered by
an order of magnitude because the original grid diverged at every point on
synthetic data. The correction is dated in the specification and appears in
the manuscript's deviations ledger.

\subsection{Uncertainty mechanisms in full}
\label{supp:methods-mechanisms}

Each mechanism turns a fitted family into predictive sample paths on the $m_x$
scale.

\begin{enumerate}
\item \textbf{Model-native.} The predictive distribution the reference paper
  itself supplies: random-walk innovation and drift uncertainty for the
  Lee--Carter-type models, the VAR's coefficient and innovation draws, the GP
  posterior, and the count distribution of the distributional head.
\item \textbf{Semiparametric Poisson bootstrap.} Refit the family $B=200$
  times to pseudo-deaths $D^{(b)}\sim\mathrm{Poisson}(E\hat m)$ and pool the
  refits' native samples \citep{brouhns2005bootstrapping}, adding the
  parameter uncertainty the native mechanism conditions away.
\item \textbf{Deep ensemble.} $M=10$ independently seeded trainings with
  identical hyperparameters, paths distributed over members with equal weight
  and drawn from each member's own predictive
  \citep{lakshminarayanan2017simple}: a mixture of near-deltas for the point
  families, a mixture of NB laws for the distributional head. Stated in
  advance: such mixtures are narrow, and if they under-cover that is a finding
  about the mechanism, not a defect to widen away.
\item \textbf{Monte Carlo dropout.} One training, with dropout left active at
  prediction time and stochastic forward passes used as paths
  \citep{gal2016dropout}, \emph{plus} the family's own innovation noise where
  it has any---so that for the LSTM, for instance, the mechanism adds
  weight uncertainty to the innovation noise rather than replacing it.
\item \textbf{Split conformal.} Absolute-residual conformal intervals on the
  log-rate scale \citep{lei2018distribution}, Mondrian-stratified by age band
  (0--24, 25--64, 65--99) and by horizon \citep{vovk2005algorithmic}. The
  mechanics: the last nine training years are held out; the base family is fit
  on the years before them; its median forecast at horizons $1,\dots,9$ is
  scored against the held-out years; the finite-sample quantile of the
  residuals gives a radius per (horizon, band); and that radius is applied
  around the median forecast of a refit on the \emph{full} training panel, so
  that the reported interval is centred on a model that has seen every
  training year.
\item \textbf{EnbPI.} Ensemble batch prediction intervals
  \citep{xu2021conformal}, adapted to annual multi-step forecasting with four
  deviations, each documented because each is a departure from the source
  algorithm:
  \begin{enumerate}
  \item $K=10$ members are fit on \emph{nested trailing-block} splits with
    origins one year apart, rather than on interior leave-block-out splits;
    interior blocks would put holes in a panel that random-walk forecasters
    cannot fit.
  \item Residuals are horizon-matched rather than pooled across steps, because
    a five-year-ahead residual and a one-year-ahead residual are not
    exchangeable in a trending panel.
  \item The interval is centred on a full-panel refit rather than on an
    aggregation of the stale members.
  \item There is no online residual update, because the test window is scored
    exactly once.
  \end{enumerate}
  What the member residuals supply, and split conformal lacks, is ten
  realisations of the latent period-effect path per horizon.
\item \textbf{Copula / joint-path conformal.} A simultaneous band over
  horizons from a scaled supremum-norm nonconformity score on those same
  member residuals \citep{diquigiovanni2022conformal}, in the spirit of
  \citet{sun2024copula} and \citet{messoudi2021copula}: each horizon's
  residual is divided by that horizon's own scale before the supremum is
  taken, so that the band widens with horizon instead of being dominated by
  the longest one, and the calibrated radius is then applied on the same
  scale. It is the only mechanism whose band is calibrated to be joint over
  $h=1,\dots,H$.
\end{enumerate}

\paragraph{The four conformal disciplines.} Calibration residuals are drawn
only from years at or before the training cutoff, on an inner time split in
which the calibration years are later than the years the base fit is trained
on, so no mechanism sees a test year. The calibration horizon is nine in every
regime, at least the longest scored horizon, so no placebo horizon reuses a
radius calibrated for a shorter one (Addendum~3~\S9). Calibration residuals
use the half-count convention of the scorer. And a wrapper yields an interval,
not a distribution: coverage, width, interval score and joint path coverage
are read from its own bounds at its own construction level of 95\%, with no
Poisson composition and no quantile-of-samples step (Addendum~2~\S3,
Addendum~3~\S6), and the 50\% and 80\% columns are not applicable. To preserve
the single sample interface for everything else, the harness draws uniformly
inside each cell's interval, independently across cells; the CRPS, log score
and PIT computed from those draws describe a placeholder density, are flagged
\texttt{secondary}, and are never ranked against distributional mechanisms
(Supplementary Section~S5).

\subsection{The crossed grid: floors and cost}
\label{supp:methods-grid}

\paragraph{Design floors.} A cell's construction implies a minimum panel
length, and the floor is a property of the family--mechanism \emph{pair},
because a wrapper's inner fits inherit the base family's own minimum. Every
family needs at least ten training years; the banded VAR needs more than ten,
the CNN more than fifteen (its ten-year window plus the five inner-validation
years) and the LSTM more than seventeen (window, inner validation and two
more), under every mechanism. Split conformal fits the base family on the
$T-9$ years beneath its nine calibration years; EnbPI and the copula band fit
their earliest staggered member on $T-18$ years. The GP's forty-year floor is
much the highest and is an identification constraint, as
Section~\ref{supp:methods-families} explains. Which cells these floors remove,
population by population and mechanism by mechanism, is tabulated in
Supplementary Section~S6.

\paragraph{Compute cost.} The three conformal columns do not reuse the native
fit. Every wrapper refits the base family---twice for split conformal, eleven
times for EnbPI and eleven for the copula band---so the honest cost of a
classical family at one origin is roughly $1+201+2+11+11$ fits: trivial for
the classical families, and the same multiplier on a neural base, where the
member refits rather than the $M=10$ ensemble dominate. The expensive axis is
neural families $\times$ members $\times$ expanding origins, which is why the
sweep is organised by origin rather than by family.

\subsection{One interface, one scoring path}
\label{supp:methods-interface}

Every family--mechanism cell implements the same two calls,
\[
  \texttt{fit}(D, E) \quad\text{and}\quad
  \texttt{sample\_mx}(h, n, \text{rng}) \;\to\; [n, h, n_{\text{ages}}],
\]
returning $n$ pathwise trajectories of $m_{x,T+1},\dots,m_{x,T+h}$. Paths, not
marginals: joint path coverage is only computable because a cell returns
trajectories.

The scoring module then makes every distributional cell
\emph{Poisson-inclusive}: on each path it draws
$D^{\ast(j)}_{x,h}\sim\mathrm{Poisson}(E_{x,T+h}\,m^{(j)}_{x,T+h})$ and forms
the half-count log crude rate, so that the object scored is the predictive law
of the rate that is actually observed (Addendum~2~\S1). The rationale is
measurable rather than aesthetic. Scoring a latent-rate sample against an
observed crude rate counts Poisson observation noise as miscalibration, most
severely at young ages and in small populations where expected deaths per cell
are small; on the harness's synthetic truth, a correctly specified Poisson
Lee--Carter covered a fraction of nominal at every level before this was
fixed, and mechanisms calibrated on observed residuals ranked better for a
reason that had nothing to do with any shift. The mirror-image error is
scoring fractional observed deaths against an integer lattice, which
Addendum~3~\S5 measures as moving the PIT KS statistic from 0.019 to 0.133 at
$\lambda=2$. Both sides of every rate-scale score therefore share one lattice
and one continuity correction.

Point forecasts, intervals, proper scores, PIT values, joint path coverage and
the life-table functionals are all computed from these samples by one
evaluation module. No model has a bespoke metric implementation, so no model
can be accidentally flattered by one. One point functional serves every point
metric: the pointwise median of the predictive samples---of the log rate per
(horizon, age) cell for RMSE and MAE, and of the per-sample life-table
functionals for $e_0$, $e_{65}$ and $\annuity$. The median is invariant under
monotone transforms, lies inside the reported interval by construction, is the
centring convention of the conformal wrappers, and is unaffected by the heavy
right tail of exponentiated Gaussian paths that would drag a sample mean. No
column uses the mean.

\subsection{Seeds and provenance}
\label{supp:methods-seeds}

The global seed is 20260825. Each cell derives its own generator from a seed
sequence over (global seed, origin, population, sex, family, mechanism), and
ensemble member $j$ is seeded from the $j$th spawn of that cell's sequence.
This is the registration's ``global seed plus member index'' in substance---
member seeds are distinct, deterministic given the global seed, and recorded
per member---while avoiding the seed collisions an additive offset produces
between neighbouring cells, and it lets any single cell be re-run in isolation
and reproduce its row. The departure from the literal registered sentence is
listed in the manuscript's deviations ledger.

Written to the results file alongside the metrics, for every cell: every seed,
the device the cell ran on, the number of ages and the number of (horizon,
age) cells scored, the training window actually used, and the exception string
of any cell that failed. Those columns are what make the error classes of
Supplementary Section~S6 auditable rather than asserted, and what allows
the model-implied joint coverage of \Hthree\ to be recomputed from the same
paths later.

\subsection{Estimators and metric derivations}
\label{supp:methods-metrics}

\paragraph{Point accuracy.} With $\hat y_{x,h} = \hat F^{-1}_{x,h}(\tfrac12)$
the pointwise predictive median,
\begin{equation}
  \RMSE = \Big(\overline{(\hat y_{x,h}-y_{x,h})^2}\Big)^{1/2},\qquad
  \MAE = \overline{|\hat y_{x,h}-y_{x,h}|},
\end{equation}
together with the error in period life expectancy at birth, $\hat e_0-e_0$.
MAE is defined and emitted for completeness beside RMSE; no result in the
manuscript is stated in terms of it, and the load-bearing fact about it is
that CRPS reduces to MAE for a point forecast, which is what makes the
RMSE-versus-CRPS comparison of \Hone\ well posed.

\paragraph{CRPS.} For a predictive distribution $F$ and outcome $y$ the
threshold and kernel representations are
\begin{equation}
  \CRPS(F,y) = \int_{-\infty}^{\infty}\big(F(u)-\ind\{u\ge y\}\big)^2\,du
  = \E_F|Z-y| - \tfrac12\,\E_F|Z-Z'|,
\end{equation}
with $Z,Z'\sim F$ independent \citep{gneiting2007strictly}. From samples we use
the plug-in estimator with sorted samples $z_{(1)}\le\dots\le z_{(n)}$,
\begin{equation}
  \widehat{\CRPS} = \frac1n\sum_{j=1}^{n}|z^{(j)}-y|
  - \frac{1}{n^2}\sum_{j=1}^{n}\big(2j-n-1\big)\,z_{(j)},
  \label{eq:crps-sample}
\end{equation}
in which the second term is the $O(n\log n)$ form of the pairwise-difference
expectation: sorting once replaces the $O(n^2)$ double sum. CRPS is computed
on the rate scale from the Poisson-inclusive samples for every distributional
cell, and on the count scale as a registered sensitivity---the same estimator
applied to the count samples $D^{\ast(j)}_{x,h}$ against the \emph{unrounded}
$D_{x,T+h}$, which needs no rounding convention at all.

\paragraph{Randomised PIT.} A forecaster is probabilistically calibrated if
$F(Y)$ is uniform \citep{gneiting2007probabilistic,dawid1984present}. Because
the predictive is an empirical distribution of $n$ samples on a lattice and
the target may tie with them, we use the randomised PIT of
\citet{czado2009predictive},
\begin{equation}
  u_{x,h} = \frac{\#\{j: z^{(j)}_{x,h} < y_{x,h}\} + V\,\big(\#\{j: z^{(j)}_{x,h}=y_{x,h}\}+1\big)}{n+1},
  \qquad V\sim\mathrm{Uniform}(0,1),
\end{equation}
which is exactly uniform under the null for any finite $n$: the randomisation
$V$ spreads the probability mass that a tie would otherwise pile on a single
lattice point, so that no part of the observed non-uniformity can be an
artefact of discreteness or of a finite sample size. Reported per cell are the
ten-bin histogram, the Kolmogorov--Smirnov distance from uniformity, and the
shape---U (over-confident), hump (under-confident), slope (biased). The KS
$p$-value is descriptive only (Addendum~2~\S4).

\paragraph{Murphy decomposition.} For a proper score $S$,
$\overline{S} = \mathrm{UNC} - \mathrm{RES} + \mathrm{MCB}$
\citep{murphy1973new,pohle2020murphy}. Two decompositions are emitted per
cell. The first is Murphy's original partition of the Brier score, taking as
outcome the 95\% interval-hit indicator and as forecast the constant nominal
level 0.95 that every cell asserts; with a constant forecast a single bin is
occupied, the three-term identity is exact, and the partition degenerates to
$\mathrm{MCB}=(0.95-\text{coverage})^2$ with $\mathrm{RES}=0$, so it restates
the coverage gap on the Brier scale and is reported for that reason only. The
second is the divergence form \citep{broecker2009reliability} on the PIT scale
with the uniform distribution as reference: the ten-bin PIT histogram is
scored against the flat claim every calibrated forecaster makes, and the
emitted miscalibration term is the squared distance of the histogram from
uniformity---the departure a single level's coverage gap cannot see. Under
that reference the reliability and resolution terms coincide numerically,
because the PIT transform standardises every forecaster to the same claim and
discards sharpness. No resolution component with independent content is
emitted: one would require a sample-based partition such as the rank-histogram
CRPS decomposition of \citet{hersbach2000decomposition}, and the isotonic
(CORP) estimator of \citet{dimitriadis2021stable} is the stable alternative
and is likewise not emitted. Both decompositions are descriptive, and no
hypothesis rests on either.

\paragraph{Wild cluster bootstrap weights.} The null distribution of the
Diebold--Mariano statistic is obtained by the wild cluster bootstrap of
\citet{cameron2008bootstrap} with the null imposed, cluster-level weights
drawn from the six-point distribution of \citet{webb2023reworking} and
$B=4999$ replications. Six-point rather than Rademacher weights because at
$G=20$ clusters, and fewer wherever the admissibility floor or the common-cell
rule removes a population, Rademacher weights supply too few distinct sign
patterns to resolve the tails of the bootstrap distribution
\citep{mackinnon2017wild}.

\paragraph{Model confidence set.} Following \citet{hansen2011model}, the set at
confidence $1-\alpha = 0.90$ is built from the $T_{\max}$ statistic with the
$e_{\max}$ elimination rule, the bootstrap distribution of the statistic being
obtained by resampling \emph{populations} as blocks (2{,}000 replications),
the analogue of their stationary block bootstrap for this panel. Sets are
formed within a common age support only---the full-age classical families,
with a second set admitting CBD only where its restricted support does not
make the losses incomparable---and on the loss appropriate to the arms:
rate-scale CRPS among distributional arms, and the 95\% interval score among
the three conformal mechanisms within each family. The
native-versus-split-conformal contrast within each family is a
Diebold--Mariano test on the interval score.

\subsection{Actuarial functionals: what is integrated, and from which paths}
\label{supp:methods-actuarial}

The life-table functionals are computed from the \emph{latent} rate paths
$m^{(j)}_{\cdot,T+1}$ of each distributional cell, before the Poisson
composition that the rate-scale scores use. The reason is that a
Poisson-inclusive draw describes one year's realised crude rate in one
population, whereas a life table is a mortality \emph{schedule}: composing
observation noise into it would price a year of sampling variation into a
reserve. The realised values are computed from the observed $m_{x,T+1}$ by the
same routine and the observed-table convention of Addendum~3~\S5, so any
life-table approximation cancels in the comparison, and both sides use the
maximal contiguous scored age range from age~0, with the range recorded.

Conformal cells have no such paths---their samples are uniform draws inside a
band, a convention for interval metrics only---so their derived-quantity rows
are reported as not applicable rather than tabulated as a distributional
claim, and every count in the manuscript's actuarial section is over the
twenty distributional arms.

\Hfive\ is restricted by Addendum~3~\S8 to the ex-post coverage of these
derived intervals, with \citet{cairns2011mortality}'s attenuation
finding---that life-table integration averages a rate-surface failure
away---as the named competitor. Both outcomes are informative, and it is worth
saying which is which: if the derived intervals miscover at least as badly as
the rate surface, the reserve is wrong; if the integration damps the failure,
the audit's headline is a statement about rates rather than about liabilities.
